\documentclass[12pt,a4paper]{report}
\usepackage{graphicx}


\begin{document}

\newcommand*{\getUniversity}{Technische Universität München}
\newcommand*{\getFaculty}{Fraunhofer Institut (AISEC)}
\newcommand*{\getTitle}{ZX-Calculus}
\newcommand*{\getAuthor}{Manuel Lerchner}
\newcommand*{\getDoctype}{Seminararbeit}
\newcommand*{\getSubmissionDate}{\today}

\begin{titlepage}
    \centering

    \includegraphics[height=20mm]{../logos/tum.pdf}
    \vspace*{20mm}

    \vspace{5mm}
    {\huge\MakeUppercase{\getFaculty{}}}\\

    \vspace{5mm}
    {\large\MakeUppercase{\getUniversity{}}}\\

    \vspace{20mm}
    {\Large \getDoctype{}}

    \vspace{15mm}
    {\huge\bfseries \getTitle{}}

    \vspace{15mm}
    \begin{tabular}{l l}
        Author:          & \getAuthor{}         \\
        Submission Date: & \getSubmissionDate{} \\
    \end{tabular}


    \vfill

    \includegraphics[height=20mm]{../logos/faculty.png}
\end{titlepage}

\end{document}